\[ %%%%%%%%%%%%%%%%%%%%%%%%%%%%%%%%%%%%%%%%%%%%%%%%%%%%%%%%%%%%%%%%%%%%%%%%%%%%%%%% \subsection{Training Configuration} \label{sec:training_configuration} %%%%%%%%%%%%%%%%%%%%%%%%%%%%%%%%%%%%%%%%%%%%%%%%%%%%%%%%%%%%%%%%%%%%%%%%%%%%%%%% This section describes the training configuration used for evaluating PHYSGEN and all baseline models. The objective is to provide a reproducible experimental protocol and ensure that performance differences are primarily determined by model capabilities rather than optimization settings. %%%%%%%%%%%%%%%%%%%%%%%%%%%%%%%%%%%%%%%%%%%%%%%%%%%%%%%%%%%%%%%%%%%%%%%%%%%%%%%% \subsubsection{Implementation Environment} %%%%%%%%%%%%%%%%%%%%%%%%%%%%%%%%%%%%%%%%%%%%%%%%%%%%%%%%%%%%%%%%%%%%%%%%%%%%%%%% All experiments are implemented using the PyTorch deep learning framework. The computational environment consists of: \begin{itemize} \item Python-based scientific computing stack; \item PyTorch for neural network optimization; \item PyTorch Geometric for graph-based operations; \item CUDA acceleration for GPU training. \end{itemize} The same computational environment is used for PHYSGEN and baseline models. %%%%%%%%%%%%%%%%%%%%%%%%%%%%%%%%%%%%%%%%%%%%%%%%%%%%%%%%%%%%%%%%%%%%%%%%%%%%%%%% \subsubsection{Optimization Strategy} %%%%%%%%%%%%%%%%%%%%%%%%%%%%%%%%%%%%%%%%%%%%%%%%%%%%%%%%%%%%%%%%%%%%%%%%%%%%%%%% PHYSGEN parameters are optimized by minimizing the complete objective function: \begin{equation} \mathcal{L}_{PHYSGEN} = \mathcal{L}_{data} + \lambda_p\mathcal{L}_{physics} + \lambda_g\mathcal{L}_{graph} + \lambda_s\mathcal{L}_{stability} + \lambda_r\mathcal{L}_{rollout}. \end{equation} The Adam optimizer is used: \begin{equation} \theta_{t+1} = \theta_t - \eta \nabla_\theta \mathcal{L}. \end{equation} The initial learning rate is: \begin{equation} \eta=10^{-3}. \end{equation} A cosine learning rate scheduler is applied to improve convergence: \begin{equation} \eta_t = \eta_0 \frac{1+\cos(\pi t/T)} {2}. \end{equation} %%%%%%%%%%%%%%%%%%%%%%%%%%%%%%%%%%%%%%%%%%%%%%%%%%%%%%%%%%%%%%%%%%%%%%%%%%%%%%%% \subsubsection{Network Architecture Parameters} %%%%%%%%%%%%%%%%%%%%%%%%%%%%%%%%%%%%%%%%%%%%%%%%%%%%%%%%%%%%%%%%%%%%%%%%%%%%%%%% The default PHYSGEN configuration is: \begin{itemize} \item graph encoder dimension: $d=128$; \item number of message passing layers: $L=6$; \item latent physical state dimension: $d_z=128$; \item activation function: GELU; \item normalization: Layer Normalization. \end{itemize} The general graph evolution block is: \begin{equation} H^{l+1} = MP_{\theta}^{l}(H^l,E). \end{equation} %%%%%%%%%%%%%%%%%%%%%%%%%%%%%%%%%%%%%%%%%%%%%%%%%%%%%%%%%%%%%%%%%%%%%%%%%%%%%%%% \subsubsection{Training Schedule} %%%%%%%%%%%%%%%%%%%%%%%%%%%%%%%%%%%%%%%%%%%%%%%%%%%%%%%%%%%%%%%%%%%%%%%%%%%%%%%% The models are trained for: \begin{equation} N_{epochs}=500 \end{equation} epochs with early stopping based on validation loss. The validation criterion is: \begin{equation} \mathcal{L}_{val} = \mathcal{L}_{PHYSGEN} ( \mathcal{D}_{val} ). \end{equation} Training is terminated if validation performance does not improve for a fixed patience interval. %%%%%%%%%%%%%%%%%%%%%%%%%%%%%%%%%%%%%%%%%%%%%%%%%%%%%%%%%%%%%%%%%%%%%%%%%%%%%%%% \subsubsection{Batching Strategy} %%%%%%%%%%%%%%%%%%%%%%%%%%%%%%%%%%%%%%%%%%%%%%%%%%%%%%%%%%%%%%%%%%%%%%%%%%%%%%%% Since physical graphs may have different sizes, mini-batches are created using graph-level batching. A batch is represented as: \begin{equation} B = \{ G_1,G_2,...,G_n \}. \end{equation} The default batch size is selected according to graph complexity: \begin{itemize} \item small ODE systems: 64 trajectories; \item medium PDE systems: 16--32 graph samples; \item large fluid simulations: 4--8 graph samples. \end{itemize} %%%%%%%%%%%%%%%%%%%%%%%%%%%%%%%%%%%%%%%%%%%%%%%%%%%%%%%%%%%%%%%%%%%%%%%%%%%%%%%% \subsubsection{Loss Weight Configuration} %%%%%%%%%%%%%%%%%%%%%%%%%%%%%%%%%%%%%%%%%%%%%%%%%%%%%%%%%%%%%%%%%%%%%%%%%%%%%%%% The initial loss coefficients are selected as: \begin{equation} \lambda_p=1.0, \end{equation} \begin{equation} \lambda_g=0.5, \end{equation} \begin{equation} \lambda_s=1.0, \end{equation} \begin{equation} \lambda_r=0.5. \end{equation} During training, adaptive weighting may modify these values according to the current dynamical regime. %%%%%%%%%%%%%%%%%%%%%%%%%%%%%%%%%%%%%%%%%%%%%%%%%%%%%%%%%%%%%%%%%%%%%%%%%%%%%%%% \subsubsection{Long-Horizon Training Strategy} %%%%%%%%%%%%%%%%%%%%%%%%%%%%%%%%%%%%%%%%%%%%%%%%%%%%%%%%%%%%%%%%%%%%%%%%%%%%%%%% To optimize long-term forecasting capability, the prediction horizon is gradually increased during training. The curriculum strategy is: \begin{equation} K_1<K_2<...<K_n, \end{equation} where $K$ represents the rollout horizon. Early training focuses on short-term prediction: \begin{equation} K=5, \end{equation} while final training evaluates extended horizons: \begin{equation} K=50-100. \end{equation} This approach improves stability during autonomous rollout. %%%%%%%%%%%%%%%%%%%%%%%%%%%%%%%%%%%%%%%%%%%%%%%%%%%%%%%%%%%%%%%%%%%%%%%%%%%%%%%% \subsubsection{Reproducibility Settings} %%%%%%%%%%%%%%%%%%%%%%%%%%%%%%%%%%%%%%%%%%%%%%%%%%%%%%%%%%%%%%%%%%%%%%%%%%%%%%%% To ensure reproducibility, all experiments use: \begin{itemize} \item fixed random seeds; \item deterministic data generation; \item documented hyperparameters; \item identical evaluation protocols. \end{itemize} The complete configuration is released together with the experimental code. %%%%%%%%%%%%%%%%%%%%%%%%%%%%%%%%%%%%%%%%%%%%%%%%%%%%%%%%%%%%%%%%%%%%%%%%%%%%%%%% \subsubsection{Summary} %%%%%%%%%%%%%%%%%%%%%%%%%%%%%%%%%%%%%%%%%%%%%%%%%%%%%%%%%%%%%%%%%%%%%%%%%%%%%%%% The training configuration defines a reproducible optimization protocol for PHYSGEN. By combining graph-based learning, physics-guided objectives, stability regularization, and long-horizon curriculum training, the experimental setup is designed to evaluate the true capabilities of physics-guided dynamical system learning. \] 
